\chapter{Theoretical Framework and Literature Review}
\label{ch:background}

This chapter introduces the concepts and the prior work on which the thesis builds. It starts from the
role of traffic signals in automated driving and the rules that give them meaning
(\cref{sec:bg-signals}). It then reviews traffic light recognition, from hand-crafted pipelines to
deep detectors (\cref{sec:bg-tlr}), and the central topic of this work: how a system decides
\emph{which} lights are relevant for the ego vehicle and which maneuvers they control
(\cref{sec:bg-relevance}). The detection and backbone architectures used later are covered in
\cref{sec:bg-detection,sec:bg-backbones}, the public datasets in \cref{sec:bg-datasets}, and the
chapter closes by positioning this work with respect to the literature (\cref{sec:bg-position}).

\section{Traffic signals in automated driving}
\label{sec:bg-signals}

Automated driving systems are commonly classified by the SAE J3016 levels of driving
automation~\cite{sae2021j3016}. From level~3 upwards the system, not the human, is responsible for
monitoring the environment within its operational design domain, so an urban system must interpret
traffic signals on its own. Vehicle-to-everything (V2X) communication could in principle broadcast
signal phases directly, but its deployment is still sparse, which leaves the camera as the only widely
available sensor for this task~\cite{polley2024tldready,polley2025atlas}. Unlike most other road
users, traffic signals carry no physical threat: their meaning is entirely conventional, defined by
national regulations that build on the 1968 Vienna Convention on Road Signs and
Signals~\cite{unece1968vienna}.

The conventions relevant to this thesis are those that link a lamp to a maneuver:
\begin{itemize}
  \item A \textbf{circular} lamp governs all maneuvers from the approach unless another signal
        overrides it. An \textbf{arrow} lamp governs only the maneuvers its arrow points to.
  \item In the \textbf{United States} the Manual on Uniform Traffic Control Devices
        (MUTCD)~\cite{fhwa2009mutcd} allows, in most jurisdictions, a right turn on a red circular
        signal after a full stop, unless a sign or a dedicated red right arrow prohibits it. A red
        circle therefore does not by itself stop a right turn.
  \item In \textbf{Germany} the Stra\ss enverkehrs-Ordnung (StVO)~\cite{stvo} makes the right turn
        obey the circular signal; turning right on red is allowed only where a green arrow sign
        (\emph{Gr\"unpfeil}) is mounted next to the light, and even then only after stopping.
  \item A \textbf{U-turn}, where legal, usually follows the left-turn signal; dedicated U-turn arrows
        are rare.
  \item The meaning of \textbf{yellow} differs: in Germany a red--yellow phase announces green, while in
        the United States yellow announces red.
\end{itemize}
The physical appearance of signals differs as well. European heads are mostly vertical three-lamp
housings mounted before the stop line, whereas US heads are often mounted on the far side of the
intersection, come in vertical and horizontal variants and are frequently seen from the back when
they serve the opposite approach~\cite{trinci2026color}. These differences are not cosmetic. A model
that learns ``right is free unless a right arrow is red'' or ``the largest light is the relevant one''
from data of one country will make wrong decisions in another. \Cref{ch:problem} returns to this point
when the label semantics are defined.

\section{Traffic light recognition}
\label{sec:bg-tlr}

Traffic light recognition (TLR) traditionally comprises the \emph{localization} of the lights in the
image and the \emph{classification} of their state (red, yellow, green, off) and, where present, of
their pictogram (circle or arrow direction).

\paragraph{Classical pipelines.} Early systems detected traffic lights with color segmentation in
suitable color spaces, followed by shape and blob analysis and temporal filtering. The survey of Jensen
et al.~\cite{jensen2016vision} reviews these approaches and identifies the problems that remain
relevant today: lights are small, their appearance depends strongly on exposure and on the camera, and
other light sources (tail lights, street lamps, reflections) look similar.

\paragraph{Deep detection.} With convolutional detectors, traffic lights became a standard, if
difficult, object detection problem; the survey of Pavlitska et al.~\cite{pavlitska2023survey} groups
the methods into adapted generic detectors, multi-stage pipelines (a detector followed by a state
classifier) and task-specific single-stage designs. Behrendt et al.~\cite{behrendt2017deep}
introduced the Bosch Small Traffic Lights Dataset together with a detection, tracking and
classification pipeline; Müller and Dietmayer~\cite{muller2018detecting} adapted single-shot detection
to the extreme aspect ratios and small sizes of lights.

\paragraph{Benchmarking across datasets.} Polley et al.\ (TLD-READY)~\cite{polley2024tldready} trained
YOLOv7, YOLOv8 and RT-DETR detectors on four datasets (BSTLD, LISA, DTLD and a private Karlsruhe set)
with a common protocol that evaluates on data unseen during training and keeps small and side-facing
lights in the test sets. For DTLD they extended YOLOv8 with the stride-4 feature level, to improve
small-light recall, and with two separate heads for state and pictogram. Their deployment on a research
vehicle exposed two weaknesses of public data: a model trained on DTLD reached only 0.28 mAP on frames
from the deployment camera (0.96 after fine-tuning on 9,000 annotated frames), and pictograms that no
public dataset contains, such as straight-right arrows, could not be recognized at all.

\paragraph{ATLAS.} The same group addressed these gaps with the ATLAS dataset and a complete
perception framework~\cite{polley2025atlas}. ATLAS annotates all vehicle lights in three time-synchronized
cameras of different focal length, covers every state--pictogram combination found in Germany
(including red--yellow and straight-right) and includes rain. A YOLO11 detector trained on DTLD and
ATLAS together keeps its DTLD accuracy while raising mAP@50 on the ATLAS test split from about 0.53 for
DTLD-only training to about 0.7, a direct measure of how far a single-camera dataset is from a
deployment setup. In the vehicle, one detector processes the medium camera and a second one alternates
between the tele camera (far lights) and the wide-angle camera (lights overhead at the stop line). The
framework is discussed further in \cref{sec:bg-maps}, since its relevance decisions come from an HD
map.

All these systems output light boxes with colors and pictograms. They leave open the question that the
planner actually needs answered: which of the lights matter for the ego vehicle, and for which
maneuver.

\section{Relevance: from lights to lanes and maneuvers}
\label{sec:bg-relevance}

A complex intersection can show a dozen lights at once, controlling different lanes, the opposite
approach or the next junction, and they frequently show different colors. Relevance estimation
decides which of them govern the ego vehicle. The literature offers four families of solutions, which
differ in the information they require at inference time; \cref{tab:relevance-methods} summarizes them.

\begin{table}[t]
\centering
\caption{Approaches to traffic light relevance. ``Inference input'' is what the method needs while
driving; ``Output'' is the granularity of the answer.}
\label{tab:relevance-methods}
\footnotesize
\begin{tabularx}{\textwidth}{@{}>{\raggedright\arraybackslash}p{3.0cm}>{\raggedright\arraybackslash}X>{\raggedright\arraybackslash}X>{\raggedright\arraybackslash}p{2.4cm}@{}}
\toprule
Method & Inference input & Output & Data \\
\midrule
Fairfield, Urmson~\cite{fairfield2011traffic} & image, prior 3D light map, localization & state of mapped lights & own \\
Possatti et al.~\cite{possatti2019traffic} & image, route-specific light map & stop/go of the route's lights & own \\
ATLAS framework~\cite{polley2025atlas} & two camera streams, HD map with signal groups & state per signal group & ATLAS \\
Langenberg et al.~\cite{langenberg2019deep} & image + lights, lane lines, lane arrows (annotated) & lights assigned to ego and neighbor lanes & DTLD + own labels \\
TLD-READY~\cite{polley2024tldready} & image (detected lights and road arrows) & relevant lights via arrow--pictogram matching & DTLD + own labels \\
Trinci et al.~\cite{trinci2026color} & image & per-light relevance and state; image-level stop/go/none & VZC-TLD, DTLD \\
TopoNet, TopoMLP~\cite{li2023graph,wu2024topomlp} & multi-camera images & lane graph and lane--light links & OpenLane-V2 \\
Zou et al.~\cite{zou2026oneperception} & front image & color and stop/go per maneuver & OpenLane-V2 \\
\bottomrule
\end{tabularx}
\end{table}

\subsection{Map-based relevance}
\label{sec:bg-maps}

Production-oriented systems sidestep most of the problem with prior maps. Fairfield and
Urmson~\cite{fairfield2011traffic} triangulate the 3D position of every light from mapping drives and
project it into the image during driving, enlarged to absorb localization errors, so that recognition
reduces to classifying the state of a small image region. Possatti et al.~\cite{possatti2019traffic}
combine a deep detector with a LiDAR-built map that stores only the lights relevant to the mapped
routes: detections outside the projected regions are discarded. The map makes relevance trivial, but
only for the route that was driven during mapping.

The ATLAS framework~\cite{polley2025atlas} is a modern instance of this family. Each detection is
back-projected as a 3D ray and matched to the mapped light positions by a minimum-cost bipartite
assignment (Hungarian algorithm~\cite{kuhn1955hungarian}) instead of a nearest-neighbor rule, which
makes the association robust to small calibration and localization errors. Lights that control the
same movement are grouped into \emph{signal groups}; a circular buffer per light accumulates
confidence-weighted detections over time, discounting old detections and those whose pictogram does
not match the map, and the most confident light of a group sets the group state. A relevant signal
group in an unknown state is conservatively treated as red. In 230 minutes of fully automated driving
in Karlsruhe the framework required no intervention, reached 99.3\% state accuracy within 120\,m and
reacted to phase changes in 184\,ms on average.

Maps solve both detection and relevance where they are accurate, but they must be built for every
intersection, kept current and localized against, and they fail silently when the infrastructure
changes. The remaining families try to infer relevance from the image itself.

\subsection{Map-free relevance estimation}

\paragraph{Heuristics.} The simplest map-free rules declare the largest, the topmost or the most
central light relevant. They are attractive because they need no training, but they fail whenever an
overhead light of another lane or the back of a far-side signal is more prominent than the ego lane's
light~\cite{polley2025atlas,trinci2026color}. Trinci et al.\ show that the largest-light rule works
reasonably on DTLD, where the relevant lights are usually the largest, but poorly on US data, where the
largest visible light is often irrelevant.

\paragraph{Traffic light to lane assignment.} Langenberg et al.~\cite{langenberg2019deep} formulate
relevance as traffic light to lane assignment (TL2LA): for the ego lane and its two neighbors, a CNN
predicts an \emph{indication vector} that marks the image columns of the lights assigned to that lane.
The image is down-sampled to $256\times256$ and its road part is warped to a bird's-eye view by
inverse perspective mapping. Their key idea is \emph{deep metadata fusion}: twelve binary metadata
feature maps (the boxes of all lights and of lights with particular arrow types, lane lines, lane
arrows, lane signs) are adaptively scaled and multiplied element-wise into an early convolutional layer. On 848
annotated DTLD sequences, fusion reaches about 93\% accuracy, more than rule-based, metadata-only and
vision-only alternatives and even above a group of 40 human raters on the full test set. The
limitation, pointed out by later work~\cite{polley2024tldready,polley2025atlas}, is that the metadata
are human annotations, which are not available while driving.

\paragraph{Lane arrows as relevance cues.} TLD-READY~\cite{polley2024tldready} replaces annotated
metadata with detections. A YOLOv8 model detects the directional arrows painted on the road, which are
larger and easier to see than the lights; a gradient-boosting classifier decides from each arrow's box,
class and offset from the image center whether it belongs to the ego lane; finally, the detected light
pictograms are matched to the relevant arrows (a right arrow to a right-arrow light or, if none
exists, to the circles). The method reaches 96\% relevance accuracy on DTLD without any map, but it
needs visible lane arrows, which disappear close to the stop line or under occlusion, and its tabular
relevance classifier did not transfer to the deployment camera without new labels. Code, weights and
the arrow-relevance labels for DTLD are public\,\webcite{web:tldready}.

\paragraph{Joint detection and relevance.} Trinci et al.~\cite{trinci2026color} argue that relevance
needs scene context that a plain detector does not capture, and solve a \emph{local} and a
\emph{global} task together. On a Faster R-CNN~\cite{ren2015faster} with a ResNet-50 backbone, two
extra heads predict the relevance and the state of each detected light from its pooled features
concatenated with its box coordinates. An image-level classifier then receives the feature pyramid
enriched with \emph{detection priors}, channels that are non-zero only where lights were detected, and
outputs one of three classes: relevant red (stop), relevant green (go) or no relevant light. Gradients
from the image-level classifier are stopped before the detection heads. Compared with heuristics, a
deep regression of the relevant light's position, RT-DETRv2 with relevance-aware classes, a Faster
R-CNN with attribute heads and a zero-shot vision--language model, the joint model is the most
accurate on both VZC-TLD and DTLD; removing the image-level classifier costs 5.8 points of image-level
mAP. Two findings are directly relevant here: images with \emph{no} relevant light are the hardest
class, because the model must recognize that \emph{all} visible lights are irrelevant, and models trained
on the German DTLD transfer poorly to the US VZC-TLD and vice versa.

\paragraph{Saliency-aware training.} Other works modify only the training signal. Greer et al.\ weight
the DETR loss towards salient (relevant) lights, which improves their recall but does not output a
relevance decision~\cite{greer2023salience}; Freitas et al.\ regress the image position where the
relevant light should appear~\cite{freitas2019relevant}.

\subsection{Topology reasoning}

The OpenLane-V2 benchmark~\cite{wang2023openlanev2} turned lane--traffic-element association into a
central task of driving-scene understanding. It annotates 3D lane centerlines, traffic elements in the
front image, and two topologies: lane-to-lane connectivity and lane-to-traffic-element links. Methods
such as TopoNet~\cite{li2023graph} and TopoMLP~\cite{wu2024topomlp}, and more recently
TopoPoint~\cite{fu2026topopoint} and one-stage designs that reuse attention maps as
topology~\cite{li2025reusing}, predict both graphs jointly from the camera images. A planner could
derive the permitted maneuvers from these graphs, but only after reconstructing the lanes and their
links correctly, so errors in any intermediate prediction propagate to the decision. This thesis uses
OpenLane-V2's annotated topology not as a prediction target but as the source from which maneuver-level
labels are derived.

\subsection{Maneuver-level signal intent}

Zou et al.~\cite{zou2026oneperception} take the last step and remove all intermediate
representations. Instead of light boxes, relevance flags or a graph from which the planner would still
have to infer what is allowed, they predict, for each of the four maneuvers straight, left, right and
U-turn, a color $c\in\{\text{none},\text{red},\text{yellow},\text{green}\}$ and a passability
$p\in\{\text{stop},\text{go}\}$. The task, called \emph{directional traffic signal understanding},
is posed on a benchmark derived from OpenLane-V2 by associating the annotated lights with ego
maneuvers through the lane topology. Its direction-aware model, described in detail in
\cref{sec:method-model}, uses a Swin Transformer backbone~\cite{liu2021swin} with an FPN, a lightweight
DETR-style light detector as auxiliary task, region-of-interest tokens of the detected lights, a global
image token and four learnable direction queries that attend to these tokens. The paper also defines
the evaluation protocol adopted in this thesis (color and passability accuracy, exact match, their
mean, the Traffic Signal Score, and the False-Go and False-Stop rates; \cref{sec:prob-metrics}).

The reported comparisons are three. Among generic backbones with plain classification heads, Swin-Tiny
is the strongest (TSS 0.873); adding the direction-aware design raises the TSS to 0.880 and lowers the
False-Go rate from 0.183 to 0.137. Multimodal large language models~\cite{liu2023visual}, which have
been applied to driving as visual question answering~\cite{sima2024drivelm}, stay behind even after
fine-tuning (TSS about 0.85--0.86 for 7--11B-parameter models) and fail almost completely zero-shot.
Finally, structured pipelines are clearly weaker: TLD-READY, adapted to the four-maneuver output with a
fallback to the approach's circular signal, reaches a TSS of 0.842 with a False-Go rate of 0.302, and
topology-based methods reach 0.66--0.71. The paper concludes that direct maneuver-level prediction is
more accurate and safer than multi-stage association. Neither the benchmark labels nor the code are
public. The present thesis replicates the best configuration of this work and revisits the comparison
with structured pipelines.

\section{Object detection}
\label{sec:bg-detection}

Both models studied in this thesis contain a detector, so the main detection paradigms are summarized
here.

\paragraph{Convolutional detectors.} Two-stage detectors such as Faster R-CNN~\cite{ren2015faster}
generate region proposals and classify pooled region features; the Feature Pyramid Network
(FPN)~\cite{lin2017feature} provides a multi-scale feature hierarchy that is particularly important for
small objects, and RoIAlign~\cite{he2017mask} extracts region features without quantization. One-stage
detectors such as YOLO~\cite{redmon2016you} predict boxes densely in a single pass and dominate
deployed traffic light detection~\cite{polley2024tldready,polley2025atlas}. Both families rely on
anchors or dense priors and on non-maximum suppression (NMS) to remove duplicates.

\paragraph{Detection transformers.} DETR~\cite{carion2020end} casts detection as set prediction: a
transformer decoder~\cite{vaswani2017attention} turns a fixed number of learned object queries into
boxes and classes, and training matches predictions to ground truth one-to-one with the Hungarian
algorithm~\cite{kuhn1955hungarian}. The matching cost and the loss combine a classification term, an
$L_1$ box term and the generalized IoU~\cite{rezatofighi2019generalized}. DETR removes anchors and NMS
but converges slowly. Deformable DETR~\cite{zhu2021deformable} replaces dense attention with sparse
sampling around reference points, which speeds up convergence and helps small objects; Group
DETR~\cite{chen2023group} accelerates training with several groups of queries, each matched
one-to-one.

\paragraph{Real-time detection transformers.} RT-DETR~\cite{zhao2024detrs} and
LW-DETR~\cite{chen2024lwdetr} showed that DETR-style detectors can match or beat YOLO-type detectors at
real-time speed, although earlier traffic light studies found RT-DETR weaker than YOLO on the small
lights of DTLD~\cite{polley2024tldready} and on relevance-aware classes~\cite{trinci2026color}.
RF-DETR~\cite{robinson2025rfdetr} builds on LW-DETR with a DINOv2~\cite{oquab2024dinov2}
self-supervised vision-transformer backbone and uses weight-sharing neural architecture search to
obtain a family of models (Nano to Large) that trade accuracy for latency; it is designed to be
fine-tuned on new domains. RF-DETR Medium is the detector of the structured baseline in this thesis.

\section{Image backbones}
\label{sec:bg-backbones}

Residual networks~\cite{he2016deep} pretrained on ImageNet~\cite{deng2009imagenet} were the default
backbone for a long time. The Vision Transformer~\cite{dosovitskiy2021image} applies a transformer to
image patches, but its global attention scales quadratically with resolution, which is a problem for
the high-resolution inputs that small traffic lights require. The \textbf{Swin
Transformer}~\cite{liu2021swin} restricts attention to local windows that shift between layers and
builds a four-stage hierarchy with strides 4 to 32, so it can serve as a drop-in replacement for
convolutional backbones in detection. \textbf{ConvNeXt}~\cite{liu2022convnet} modernized the
convolutional design along the same lines (large depthwise kernels, LayerNorm, inverted bottlenecks)
and matched Swin's accuracy; \textbf{ConvNeXt V2}~\cite{woo2023convnext} adds a global response
normalization (GRN) layer and fully convolutional masked-autoencoder (FCMAE) pretraining. Swin-Tiny and
ConvNeXt V2-Tiny have nearly the same size, channels and strides, which makes them a clean pair for an
ablation (\cref{sec:res-backbone}).

\section{Datasets}
\label{sec:bg-datasets}

Public data for traffic lights fall into two groups: dedicated traffic-light datasets, which annotate
lights in detail but rarely say which ones matter, and general driving datasets, in which lights are
one category among many. \Cref{tab:datasets} summarizes the datasets considered in this work; counts
refer to annotated images and vehicle-light boxes as distributed (or, where noted, as exported to the
common per-lamp format used here). Three properties determine how useful a dataset is for
maneuver-level prediction: whether it marks the \emph{relevant} lights, whether it annotates
\emph{pictograms} (arrows), and how diverse it is in country, camera and illumination.

\begin{table}[t]
\centering
\caption{Datasets considered in this thesis. St.\ = annotated states (R red, Y yellow, G green, O off,
U unknown); Pict.\ = arrow pictograms annotated; Relev.\ = how relevance is defined; Use = role in the
experiments (``--'' = surveyed only).}
\label{tab:datasets}
\footnotesize\setlength{\tabcolsep}{3pt}
\begin{tabularx}{\textwidth}{@{}llrrlllX@{}}
\toprule
Dataset & Region & Images & Lights & St. & Pict. & Relev. & Use \\
\midrule
LISA~\cite{jensen2016vision} & San Diego, US & 36k & 109k & RYG & left & -- & OOD test \\
BSTLD~\cite{behrendt2017deep} & California, US & 13k & 24k & RYGO & train & -- & -- \\
DTLD~\cite{fregin2018dtld} & Germany & 41k & 292k$^a$ & RYG+ & yes & route & train, test \\
S$^2$TLD~\cite{yang2022scrdet} & Shanghai, CN & 5.8k & 14k & RYGO & -- & -- & -- \\
VZC-TLD~\cite{trinci2026color} & US & 3.0k & 25k & RYGU & -- & lane & -- \\
ATLAS~\cite{polley2025atlas} & Karlsruhe, DE & 33k & 73k & RYG+ & yes & -- & OOD, per light \\
BDD100K~\cite{yu2020bdd100k} & US & 80k$^b$ & 213k$^b$ & RYG & -- & -- & -- \\
OLV2 A~\cite{wang2023openlanev2,wilson2021argoverse} & US & 27k$^c$ & -- & RYG & yes & topol. & train, eval \\
OLV2 B~\cite{wang2023openlanev2,caesar2020nuscenes} & US, SG & 34k$^c$ & -- & RYG & yes & topol. & train, eval \\
\bottomrule
\multicolumn{8}{@{}p{\textwidth}@{}}{\footnotesize RYG+: also red--yellow and off (DTLD also unknown).
$^a$ all lamp types, including pedestrian, bicycle and tram lights; $^b$ train and validation splits;
$^c$ front-camera frames of the train and validation splits. OLV2 = OpenLane-V2.} \\
\end{tabularx}
\end{table}

\subsection{Traffic-light datasets}

\paragraph{LISA.} The LISA Traffic Light Dataset~\cite{jensen2016vision}\,\webcite{web:lisa} was
recorded in San Diego and consists of day and night video sequences at $1280\times960$ pixels, about
36,000 annotated frames with 109,000 lights in total. Lights carry a state (red, yellow, green) and,
for some, a left or straight arrow, but no relevance information. Consecutive frames are highly
redundant, and the boxes are known to be loose and incomplete: some visible lights are not
annotated~\cite{polley2025atlas,fregin2018dtld}. Its night sequences, which make up more than half of
the test set, make it the main test of illumination robustness in this work.

\paragraph{Bosch Small Traffic Lights (BSTLD).} BSTLD~\cite{behrendt2017deep}\,\webcite{web:bstld}
contains about 13,400 images at $1280\times720$ recorded around Palo Alto, California, with about
24,000 annotated lights, many of them very small. The training split annotates 13 state--pictogram
classes; the test split only the four states red, yellow, green and off, which limits it to color
recognition. The test split is larger than the training split (8,334 versus 5,093 images).

\paragraph{DriveU Traffic Light Dataset (DTLD).} DTLD~\cite{fregin2018dtld}\,\webcite{web:dtld} is the
largest public traffic-light dataset: about 41,000 frames from 2,110 sequences in eleven German cities,
recorded at 1\,Hz with a stereo camera at $2048\times1024$ pixels and distributed as raw Bayer
images. It annotates about 292,000 lamps of all kinds with a rich attribute set: facing direction
(front, back, left, right), state (red, yellow, red--yellow, green, off, unknown), pictogram (circle,
left, right, straight and straight-left arrows, pedestrian, bicycle, tram), number of lamps, occlusion
and a \emph{relevant} flag. The flag marks the lights that serve the route the vehicle actually drove;
the route itself is not annotated, so a light can be irrelevant only because the vehicle went another
way. All sequences come from one vehicle and camera, mostly by day and in good weather, and each
sequence approaches a single intersection. DTLD is the training set for the German-rule models of this
thesis.

\paragraph{S$^2$TLD.} The SJTU Small Traffic Light Dataset~\cite{yang2022scrdet}\,\webcite{web:s2tld}
contains 5,786 images from Shanghai at $1920\times1080$ and $1280\times720$ with about 14,000 lights in
five classes (red, yellow, green, off and ``wait on''). It adds a Chinese urban
setting with very small lights but has neither pictograms nor relevance.

\paragraph{VZC-TLD.} The Verizon Connect Traffic Light Dataset~\cite{trinci2026color}\,\webcite{web:vzc}
contains about 3,000 dashcam images ($1280\times720$, 150° diagonal field of view) taken from 1,809
commercial vehicles of different sizes across the United States, each image from a different video
and thus a different intersection. Each of its 25,000 lights is annotated with a state (red, yellow,
green, unknown) and a binary relevance flag, defined as the light that governs the lane occupied by the
ego vehicle. It is small but far more diverse than DTLD in camera height, field of view, weather and
time of day, and it is the only public US dataset with relevance labels.

\paragraph{ATLAS.} The Applied Traffic Light Annotation Set~\cite{polley2025atlas}\,\webcite{web:atlas}
contains 33,044 images with 72,998 lights from three synchronized cameras of the research vehicle CoCar
NextGen in Karlsruhe: front-medium ($61^\circ\times39^\circ$, $1920\times1200$, 25,158 images),
front-tele ($31^\circ\times20^\circ$, $1920\times1200$, 5,109 images) and front-wide
($106^\circ\times92^\circ$, $2592\times2048$, 2,777 images). Lights are annotated from 2--4 pixels of
width, corresponding to distances of up to 200\,m on the tele camera, in 25 state--pictogram classes
that include all arrow combinations, red--yellow, off and two-lamp on-demand signals; 2,776 images were
recorded in medium to heavy rain. ATLAS provides no relevance or maneuver labels, so it is used here
only for a per-light evaluation of detection, color and arrow direction across cameras.

\subsection{Driving datasets with traffic-light annotations}

\paragraph{BDD100K.} BDD100K~\cite{yu2020bdd100k}\,\webcite{web:bdd100k} consists of 100,000 driving
videos crowd-sourced in several US regions, with one annotated key frame per video ($1280\times720$;
70,000 training, 10,000 validation and 20,000 test images). Among its object categories are traffic
lights with a color attribute; the training and validation splits together contain about 213,000 of
them. It covers a wide range of weather, times of day and cameras, but lights are often tiny and it has
neither pictograms nor relevance.

\paragraph{OpenLane-V2.} OpenLane-V2~\cite{wang2023openlanev2}\,\webcite{web:openlanev2} is built on
Argoverse~2~\cite{wilson2021argoverse} (subset A: seven ring cameras at 2\,Hz in six US cities; front
camera $1550\times2048$, portrait) and nuScenes~\cite{caesar2020nuscenes} (subset B: six cameras,
scenes in Boston and Singapore; front camera $1600\times900$). Traffic elements are annotated as 2D
boxes in the front image only, with an attribute that encodes the color for circular lights and the
direction for arrows. Its value for this thesis lies in the \code{topology_lcte} (lane centerline to
traffic element) and \code{topology_lclc} (lane to lane) annotations: together with the 3D centerline
geometry they make it possible to determine which maneuver each relevant light controls, so that
maneuver-level labels can be derived without manual annotation (\cref{sec:method-labels}).

\subsection{Use in this thesis}

OpenLane-V2 is the only dataset from which maneuver-level labels can be derived for every direction,
and it is the basis of the replicated benchmark. DTLD provides the second source of relevance---the
driven route---under German rules, and is used to train and test a model in a different country.
LISA and ATLAS, which lack relevance labels, serve as out-of-distribution test sets: LISA with
maneuver labels reconstructed from light sizes and kinds (US, day and night), ATLAS per light (Germany,
three cameras). BSTLD, S$^2$TLD, VZC-TLD and BDD100K were converted to the same per-lamp format but
are not used in the experiments reported here: they either lack the pictograms needed to tell which
maneuver a light controls, or, in the case of VZC-TLD, provide relevance only at the lane level. They
are natural candidates for the multi-dataset training discussed in \cref{sec:disc-future}.

The survey also shows a gap. Only three public sources say which lights matter, and each defines
relevance differently: by the route driven (DTLD), by the lane occupied (VZC-TLD) or by the annotated
lane topology (OpenLane-V2). None provides maneuver-level labels directly, and the two datasets with
pictograms and relevance (DTLD and OpenLane-V2) follow different national traffic rules.

\section{Positioning of this work}
\label{sec:bg-position}

The literature offers, on one side, mature light detectors that do not decide
relevance~\cite{behrendt2017deep,polley2024tldready} and, on the other, systems that decide relevance
with additional infrastructure or annotation: HD maps~\cite{fairfield2011traffic,polley2025atlas},
annotated lane metadata~\cite{langenberg2019deep} or predicted lane
topology~\cite{li2023graph,wu2024topomlp}. Map-free, image-only methods have so far produced a
per-light relevance flag or a single image-level stop/go decision~\cite{polley2024tldready,trinci2026color},
which does not distinguish between maneuvers. Maneuver-level prediction from a single
image~\cite{zou2026oneperception} answers the planner's question most directly, but it has been
evaluated on a single, unreleased benchmark, against structured baselines that were designed for a
different output, and in one traffic-rule setting only.

This thesis contributes an open reconstruction of that benchmark, an independent replication, a
comparison with a structured alternative whose classes are designed for the maneuver-level output, and
the first cross-dataset evaluation of the task that we are aware of, spanning US and German traffic
rules, day and night, and three camera fields of view.
